% NeurIPS Paper Checklist

\section*{NeurIPS Paper Checklist}

\begin{enumerate}

\item {\bf Claims}
    \item[] Question: Do the main claims made in the abstract and introduction accurately reflect the paper's contributions and scope?
    \item[] Answer: \answerYes{}
    \item[] Justification: The abstract states three claims (full-distribution-aware samplers dominate judged quality at frontier scale, min-p sits in a middle tier, and the gap widens with model scale) and each is supported by specific numbers from the experimental record. The three numbered contributions in the introduction map directly to the results sections.

\item {\bf Limitations}
    \item[] Question: Does the paper discuss the limitations of the work performed by the authors?
    \item[] Answer: \answerYes{}
    \item[] Justification: Section~\ref{sec:limitations} discusses model scale range (1B--120B open-weight only, no closed frontier models), domain restriction to creative fiction, LLM-as-judge biases, the small multi-model judged subset, the quantization-only inference setup, the one-dimensional nature of the selected-token rank diagnostic, and the unfused \pless{} implementation.

\item {\bf Theory Assumptions and Proofs}
    \item[] Question: For each theoretical result, does the paper provide the full set of assumptions and a complete (correct) proof?
    \item[] Answer: \answerNA{}
    \item[] Justification: This paper is empirical and does not contain theoretical results or proofs.

\item {\bf Experimental Result Reproducibility}
    \item[] Question: Does the paper fully disclose all the information needed to reproduce the main experimental results of the paper to the extent that it affects the main claims and/or conclusions of the paper?
    \item[] Answer: \answerYes{}
    \item[] Justification: Section~\ref{sec:setup} specifies the exact models, prompt genres, output lengths, sampler configurations, evaluation metrics, and infrastructure. Section~\ref{sec:method} describes the \pless{} algorithm. All sampler parameter values are reported in the tables.

\item {\bf Open access to data and code}
    \item[] Question: Does the paper provide open access to the data and code, with sufficient instructions to faithfully reproduce the main experimental results?
    \item[] Answer: \answerYes{Code and data will be released to the public upon acceptance.}  \href{https://anonymous.4open.science/r/long_context-20AE/}{Code is available here}
    \item[] Justification: The experimental codebase and generated outputs will be released as a public repository upon acceptance.

\item {\bf Experimental Setting/Details}
    \item[] Question: Does the paper specify all the training and test details necessary to understand the results?
    \item[] Answer: \answerYes{}
    \item[] Justification: All models are used off-the-shelf with no fine-tuning. Models up to 32B are inferred with SGLang on a single RTX 5090 with 4-bit quantization. The 70B and 120B runs use FP8 quantization on a 4xH100 cluster. All sampler parameters, temperatures, output lengths, and evaluation metrics are specified in Section~\ref{sec:setup}.

\item {\bf Experiment Statistical Significance}
    \item[] Question: Does the paper report error bars suitably and correctly defined or other appropriate information about the statistical significance of the experiments?
    \item[] Answer: \answerYes{}
    \item[] Justification: We report means over multiple generations per cell (3--24 depending on the study) and note sample sizes throughout. The 16K cross-model transfer figure (Figure~\ref{fig:transfer}) reports error bars from the across-model spread. The 32K campaign averages 24 generations per cell. We discuss the small multi-model judged subset as a limitation.

\item {\bf Experiments Compute Resources}
    \item[] Question: For each experiment, does the paper provide sufficient information on the computer resources needed to reproduce the experiments?
    \item[] Answer: \answerYes{}
    \item[] Justification: Models up to 32B run on a single NVIDIA RTX 5090 (32 GB) with 4-bit quantized models served through SGLang. The 70B and 120B runs use FP8 quantization on a 4xH100 cluster. The total compute is on the order of a few hundred GPU-hours.

\item {\bf Code Of Ethics}
    \item[] Question: Does the research conducted in the paper conform, in every respect, with the NeurIPS Code of Ethics?
    \item[] Answer: \answerYes{}
    \item[] Justification: The research studies open-weight language models on creative fiction generation. No personally identifying information is collected.

\item {\bf Broader Impacts}
    \item[] Question: Does the paper discuss both potential positive societal impacts and negative societal impacts of the work?
    \item[] Answer: \answerYes{}
    \item[] Justification: Improved long-context sampling lets open-weight models produce coherent long-form text without retraining, which lowers the cost of high-quality generation for low-resource users and reduces the energy cost of post-training fixes. The same improvement applies to misuse cases (longer fluent disinformation, longer phishing). The underlying models already enable both, and our contribution is a sampler-level change that does not unlock new capabilities beyond what the base models support.

\item {\bf Safeguards}
    \item[] Question: Does the paper describe safeguards that have been put in place for responsible disclosure of data with personal or offensive content?
    \item[] Answer: \answerYes{}
    \item[] Justification: The generated outputs are creative fiction from non-adversarial standard prompts. We screened the prompt set to exclude topics that would induce graphic or unsafe content from the underlying models. The released artifact will be the sampler implementation and aggregate evaluation scripts, not the raw generations.

\item {\bf Licenses for existing assets}
    \item[] Question: Are the creators or original owners of assets used in the paper properly credited and are the license and terms of use explicitly mentioned and properly respected?
    \item[] Answer: \answerYes{}
    \item[] Justification: All models used are open-weight and properly cited. SGLang is cited. No proprietary datasets are used.

\item {\bf New Assets}
    \item[] Question: Are new assets introduced in the paper well documented and is the documentation provided alongside the assets?
    \item[] Answer: \answerYes{The \pless{} sampler implementation and evaluation scripts will be released with documentation.}
    \item[] Justification: The \pless{} algorithm is described in Section~\ref{sec:method}. Code release is planned upon acceptance.

\item {\bf Crowdsourcing and Research with Human Subjects}
    \item[] Question: For crowdsourcing experiments and research with human subjects, does the paper include the full text of instructions given to participants and screenshots, if applicable, as well as details about compensation?
    \item[] Answer: \answerNA{}
    \item[] Justification: The paper does not involve crowdsourcing or research with human subjects.

\item {\bf Institutional Review Board (IRB) Approvals or Equivalent for Research with Human Subjects}
    \item[] Question: Does the paper describe potential risks incurred by study participants, whether such risks were disclosed to the subjects, and whether Institutional Review Board (IRB) approvals were obtained?
    \item[] Answer: \answerNA{}
    \item[] Justification: The paper does not involve research with human subjects.

\item {\bf Declaration of LLM usage}
    \item[] Question: Does the paper describe the usage of LLMs if it is an important, original, or non-standard component of the core methods in this research?
    \item[] Answer: \answerYes{}
    \item[] Justification: GPT-4o is used as the LLM judge in our 21-metric evaluation protocol (Section~\ref{sec:setup}). The paper itself studies sampling algorithms for open-weight LLMs across ten models (Llama, Qwen, Mistral, GPT-OSS), all of which are documented in Section~\ref{sec:setup}. LLMs were not used in the writing of the paper beyond standard editing assistance, which does not affect the core methods or claims.

\end{enumerate}
